Una empresa fabrica dos productos químicos, $P_1$ y $P_2$, cuyos márgenes unitarios son, respectivamente, 1 y 4 euros. Semanalmente dispone de 40 toneladas de un producto crítico $C$, que es un cuello de botella. Cada tonelada de $P_1$ consume una tonelada de $C$ y cada tonelada de $P_2$ genera una tonelada de $C$.

Además por normativa, no se pueden producir más de 20 toneladas de $P_2$ y entre los dos productos se deben entregar al menos 10 toneladas.

Un experto el optimización, ha formulado el correspondiente problema de probramación lineal, donde $x_i$ representa la producción en toneladas de $P_i$. 


\begin{equation}
\begin{split}
\mbox{max. } z = x_1 + 4x_2\\
s.a.:\\
x_1 - x_2 \leq 40\\
x_1 + x_2 \geq 10\\
x_2 \leq 20\\
x_1,\,\,x_2\geq 0\\
\end{split}
\end{equation}

Además, ha identificado que el plan de producción actual de la empresa quedaría descrito por la siguiente tabla, correspondiente a una solución básica:

\begin{center}
\begin{tabular}{c|c|ccccc|}
 & $z$ & $x_1$ & $x_2$ & $h_1$ & $h_2$ & $h_3$\\ 
 & $-40$ & $-3$ & $0$ & $0$ & $4$ & $0$ \\ 
\hline
$x_2$ & 10  & $1$ & $1$ & $0$ & $-1$ & $0$\\ 
$h_1$ & 50  & $2$ & $0$ & $1$ & $-1$ & $0$\\ 
$h_3$ & 10  & $-1$ & $0$ & $0$ & $1$ & $1$\\ 
\hline
\end{tabular}
\end{center}



\begin{enumerate}
    \item \textbf{A partir de la información de la tabla, justificar que disminuir el compromiso comercial, $b_2$ en una unidad daría lugar a una disminución de la función objetivo.}
    

    El incremento (disminución) de $z$ cuando $b_2$ aumenta (disminuye) en un unidad viene dado por el precio sombra de la segunda restricción, $\pi^B_2 = V^B_{h_2}=4$ (restricción de tipo menor o igual).

    En efecto si $b_2=9$, $z$ disminuye en 4 unidades (4000 euros).


    \item \textbf{Explicar por qué relajar el compromiso comercial (lo cual implica aumentar el tamaño de la región de factibilidad) conduce a una solución con peor función objetivo.}

    La explicación de esta aparente paradoja viene dada por el hecho de que la solución no es óptima. La base a la que se refiere la tabla tiene como variables básicas: $x_2$, $h_1$ y $h_3$, es decir corresponden a la intersección de la restricción 2 con el eje de $Y$, como se aprecia en la siguiente figura.

\begin{tikzpicture}
    \begin{axis}[
        axis lines = middle,
        xlabel = {$x_1$},
        ylabel = {$x_2$},
        xmin = 0, xmax = 70,
        ymin = 0, ymax = 40,
        width=14cm, height=10cm,
        legend pos = north west,
        grid = major,
    ]
    % Restriction 1: x1 - x2 <= 40
    \addplot[name path=A, domain=0:70, thick, blue] {x - 40};
    \addlegendentry{$x_1 - x_2 \leq 40$}

    % Restriction 2: x1 + x2 >= 10
    \addplot[name path=B, domain=0:70, thick, red] {10 - x};
    \addlegendentry{$x_1 + x_2 \geq 10$}

    % Restriction 3: x2 <= 20
    \addplot[name path=C, domain=0:70, thick, green!70!black] {20};
    \addlegendentry{$x_2 \leq 20$}

    % Region Factible (shading)
    \addplot [
        thick,
        color=yellow,
        fill=yellow,
        fill opacity=0.3
    ]
    coordinates {
        (10,0) (0,10) (0,20) (60, 20) (40,0)
    };

    % Objective function line (longer)
    \addplot[dashed, thick, black, domain=-10:80] {0.25*x + 8};
    \addlegendentry{F. objetivo $z = 1x_1 + 4x_2$}

        % Adding the label for the point
    \node[circle, draw=black, minimum size=6pt, inner sep=0pt] at (axis cs:0, 10) {};
    \node at (axis cs:8, 10) [anchor=south] {Solución básica};
    
    \end{axis}
\end{tikzpicture}

    Relajar el compromiso comercial significa desplazar la restricción 2 hacia la izquierda, con lo que la solución básica, intersección de esta restricción 2 con el eje $Y$ haría que disminuyera $z$.

    La interpretación es la siguiente. En esta base:
    \begin{itemize}
        \item $x_1$ es variable no básica, con lo que se decide no producir $P_1$
        \item $h_2$ es variable no básica, con lo se está decidiendo cumplir el compromiso comercial por el mínimo acordado.
    \end{itemize}

    Si queremos producir solo lo que viene dado por el compromiso comercial ($h_2=0$) y solo se produce $P_2$ ($x_1=0$) el resultado es que se produce solo $P_2$ y en menor cantidad, con lo que la función objetivo disminuye.


    Este comportamiento no se observaría si estuviéramos en la solución óptima.



\end{enumerate}